\section{تحلیل فرکانسی و بررسی پدیده \lr{Aliasing} در کاهش رزولوشن}

در این بخش، قصد داریم کد نوشته شده برای تحلیل فرکانسی را به صورت گام‌به‌گام و خط به خط بررسی کنیم. هدف اصلی، درک نحوه استخراج پنجره فرکانس‌های مجاز و محاسبه میزان انرژی از دست رفته در فرآیند کاهش رزولوشن است.

\subsection*{۱. انتقال تصویر به فضای فرکانس}
در ابتدای حلقه، تصویر را وارد فضای فرکانس می‌کنیم:

\begin{codebox}[\lr{FFT and Shift}]
\begin{lstlisting}[language=Python]
for i, img in enumerate(images):
    print(titles[i])
    
    F = np.fft.fft2(img)
    F_shifted = np.fft.fftshift(F)
\end{lstlisting}
\end{codebox}

دستور \lr{\texttt{np.fft.fft2}} تصویر را از فضای مکانی (پیکسل‌ها) به فضای فرکانس منتقل می‌کند. در خروجی استاندارد این دستور، فرکانس صفر (بخش‌های یکنواخت و بدون تغییر تصویر) در گوشه‌های ماتریس قرار می‌گیرد. با استفاده از دستور \lr{\texttt{np.fft.fftshift}} این فرکانس صفر را دقیقاً به مرکز ماتریس منتقل می‌کنیم تا استخراج پنجره مرکزی در مراحل بعد، از نظر ریاضی و کدنویسی ساده‌تر شود.

\subsection*{۲. محاسبه انرژی کل و تعیین ابعاد}
پیش از محاسبه الیاسینگ، به ابعاد فضای فرکانس و انرژی کل تصویر نیاز داریم:

\begin{codebox}[\lr{Total Energy}]
\begin{lstlisting}[language=Python]
    M, N = F_shifted.shape
    E_total = np.sum(np.abs(F_shifted)**2)
\end{lstlisting}
\end{codebox}

متغیرهای \lr{M} و \lr{N} ابعاد ماتریس فرکانس را در خود ذخیره می‌کنند. سپس بر اساس «قضیه پارسوال»، انرژی کل سیگنال را با محاسبه مجموع توان دوم اندازه (مگنیتیود) تمامی فرکانس‌های موجود در ماتریس \lr{\texttt{F\_shifted}} به دست می‌آوریم.

\newpage
\subsection*{۳. محاسبه بازه و اندیس‌های مجاز برای نرخ ۴}
بر اساس قضیه نمونه‌برداری نایکوئیست، فرکانس‌های مجاز باید از مرکز تصویر استخراج شوند:

\begin{codebox}[\lr{Indices for Rate 4}]
\begin{lstlisting}[language=Python]
    r4 = 4
    row_start_4 = M // 2 - M // (2 * r4)
    row_end_4   = M // 2 + M // (2 * r4)
    col_start_4 = N // 2 - N // (2 * r4)
    col_end_4   = N // 2 + N // (2 * r4)
\end{lstlisting}
\end{codebox}

زمانی که تصویر را با نرخ $r=4$ کوچک می‌کنیم، طول و عرض پنجره فرکانس‌های مجاز برابر با $\frac{M}{4}$ و $\frac{N}{4}$ خواهد بود. از آنجا که مبدا مختصات ما اکنون در مرکز ماتریس (یعنی نقطه \lr{M//2} و \lr{N//2}) قرار دارد، برای رسیدن به حاشیه‌های این پنجره، باید به اندازه نصف طول آن (یعنی $\frac{M}{8}$ و $\frac{N}{8}$) از مرکز به سمت بالا، پایین، چپ و راست حرکت کنیم. این خطوط دقیقاً نقطه شروع و پایان این کادر را محاسبه می‌کنند.

\subsection*{۴. استخراج کادر اصلی و محاسبه الیاسینگ (نرخ ۴)}
اکنون با استفاده از اندیس‌های به دست آمده، مقادیر انرژی را محاسبه می‌کنیم:

\begin{codebox}[\lr{Energy Calculation Rate 4}]
\begin{lstlisting}[language=Python]
    center_box_4 = F_shifted[row_start_4:row_end_4, col_start_4:col_end_4]
    E_main_4 = np.sum(np.abs(center_box_4)**2)

    E_aliasing_4 = E_total - E_main_4
    ratio_4 = E_main_4 / E_aliasing_4
    print(f"Energy 4: {ratio_4:.4f}")
\end{lstlisting}
\end{codebox}

با استفاده از قابلیت برش (Slicing) در آرایه نامپای، ناحیه مرکزی (\lr{\texttt{center\_box\_4}}) استخراج می‌شود. انرژی این کادر مرکزی (\lr{\texttt{E\_main\_4}}) نشان‌دهنده اطلاعات مفیدی است که پس از کوچک‌سازی تصویر حفظ می‌شود. با کسر این مقدار از انرژی کل، مقدار انرژی از دست رفته (\lr{\texttt{E\_aliasing\_4}}) به دست می‌آید. این انرژی از دست رفته، روی تصویر جدید به صورت نویز و اعوجاج سایه می‌اندازد. در نهایت، نسبت انرژی حفظ شده به انرژی الیاسینگ محاسبه و چاپ می‌شود.

\newpage
\subsection*{۵. تکرار محاسبات برای نرخ ۱۶}
دقیقاً همان منطق ریاضی برای نرخ ۱۶ پیاده‌سازی می‌شود:

\begin{codebox}[\lr{Calculations for Rate 16}]
\begin{lstlisting}[language=Python]
    r16 = 16
    row_start_16 = M // 2 - M // (2 * r16)
    row_end_16   = M // 2 + M // (2 * r16)
    col_start_16 = N // 2 - N // (2 * r16)
    col_end_16   = N // 2 + N // (2 * r16)

    center_box_16 = F_shifted[row_start_16:row_end_16, col_start_16:col_end_16]
    E_main_16 = np.sum(np.abs(center_box_16)**2)

    E_aliasing_16 = E_total - E_main_16
    ratio_16 = E_main_16 / E_aliasing_16
    print(f"Energy 16: {ratio_16:.4f}")
    print("-" * 20)
\end{lstlisting}
\end{codebox}

در اینجا تنها تفاوت، اندازه پنجره مرکزی است. برای نرخ ۱۶، شعاع حرکت از مرکز معادل $\frac{M}{32}$ و $\frac{N}{32}$ خواهد بود. به همین دلیل کادر بسیار کوچکتری استخراج شده و طبیعتاً انرژی \lr{Aliasing} بسیار بیشتر (و نسبت محاسبه شده کمتر) خواهد بود.

\subsection*{تحلیل خروجی‌ها و بررسی نتایج}
با اجرای قطعه کدهای فوق روی سه تصویر ورودی، نتایج زیر حاصل شد:

\begin{itemize}
    \item \textbf{تصویر \lr{Grid}:} این تصویر کمترین نسبت انرژی را در هر دو نرخ (۸.۱۲ و ۴.۴۰) نشان می‌دهد. دلیل این امر، وجود خطوط تیز و شبکه‌ای است که معادل فرکانس‌های بسیار بالا در فضای فوریه هستند. با کاهش ابعاد تصویر، این فرکانس‌های بالا از پنجره مجاز خارج شده و بالاترین میزان الیاسینگ را ایجاد می‌کنند.
    
    \item \textbf{تصویر \lr{Fly Eye}:} در نرخ ۴، این تصویر بالاترین نسبت انرژی (۱۵۴.۱۸) را دارد؛ به این معنا که بافت چشم مگس عمدتاً از فرکانس‌های پایینی تشکیل شده که در پنجره $\frac{1}{4}$ جای می‌گیرند. اما در نرخ ۱۶، این نسبت به شدت افت کرده و به ۲۲.۸۹ می‌رسد؛ زیرا پنجره مجاز به قدری کوچک می‌شود که بخش عمده‌ای از بافت میانی تصویر از دست رفته و تبدیل به الیاسینگ می‌شود.
    
    \item \textbf{تصویر \lr{Ferris Wheel}:} این تصویر به عنوان یک منظره طبیعی، ترکیبی متعادل از نواحی یکنواخت (مانند آسمان) و لبه‌های ساختاری (مانند میله‌ها) دارد. به همین دلیل، رفتار فرکانسی آن متعادل‌تر بوده و نسبت انرژی آن بین دو تصویر دیگر قرار می‌گیرد (۱۰.۷۷ برای نرخ ۴ و ۵.۰۸ برای نرخ ۱۶).
\end{itemize}

\noindent\textbf{نتیجه‌گیری:} هرچه یک تصویر دارای لبه‌های تیزتر و بافت‌های متراکم‌تر باشد، در فضای فرکانس انرژی بیشتری در نواحی فرکانس‌بالا خواهد داشت. در نتیجه، هنگام کاهش رزولوشن، اطلاعات بیشتری از دست داده و شدت پدیده مخرب الیاسینگ در آن بیشتر خواهد بود.